\section{Code Snippets}

\begin{frame}[fragile]{Code Listing - Bash and Python}
  \begin{columns}[T]
    \begin{column}{0.48\textwidth}
      \begin{lstlisting}[language=bash, numbers=left, basicstyle=\ttfamily\tiny, title={\tiny\bfseries Shell Script Example}]
#!/bin/bash
# Deploy application to staging
set -euo pipefail

DEPLOY_ENV="${1:-staging}"
IMAGE="registry.example.com/app:latest"

echo "Deploying to ${DEPLOY_ENV}..."

podman pull "${IMAGE}"
podman run -d \
  --name myapp \
  -p 8080:8080 \
  -e ENV="${DEPLOY_ENV}" \
  "${IMAGE}"

echo "Deployment complete!"
      \end{lstlisting}
    \end{column}
    \begin{column}{0.48\textwidth}
      \begin{lstlisting}[language=python, numbers=left, basicstyle=\ttfamily\tiny, title={\tiny\bfseries Python Example}]
#!/usr/bin/env python3
"""Sample API client."""
import requests
from typing import Optional

class APIClient:
    def __init__(self, base_url: str):
        self.base_url = base_url
        self.session = requests.Session()

    def get_status(self) -> dict:
        """Fetch service status."""
        resp = self.session.get(
            f"{self.base_url}/api/v1/status"
        )
        resp.raise_for_status()
        return resp.json()

if __name__ == "__main__":
    client = APIClient("https://api.example.com")
    print(client.get_status())
      \end{lstlisting}
    \end{column}
  \end{columns}
\end{frame}

\begin{frame}[fragile]{Code Listing - YAML and JSON}
  \begin{columns}[T]
    \begin{column}{0.50\textwidth}
      \begin{lstlisting}[style=yamlstyle, basicstyle=\ttfamily\fontsize{4.5}{5.5}\selectfont, title={\tiny\bfseries Kubernetes Deployment (YAML)}]
apiVersion: apps/v1
kind: Deployment
metadata:
  name: my-application
  namespace: production
  labels:
    app: my-app
    version: v2.1.0
spec:
  replicas: 3
  selector:
    matchLabels:
      app: my-app
  template:
    metadata:
      labels:
        app: my-app
    spec:
      containers:
        - name: app
          image: registry.example.com/app:v2.1.0
          ports:
            - containerPort: 8080
          resources:
            limits:
              memory: "256Mi"
              cpu: "500m"
          livenessProbe:
            httpGet:
              path: /healthz
              port: 8080
            initialDelaySeconds: 10
      \end{lstlisting}
    \end{column}
    \begin{column}{0.50\textwidth}
      \begin{lstlisting}[language=, basicstyle=\ttfamily\fontsize{4.5}{5.5}\selectfont, title={\tiny\bfseries API Response (JSON)}]
{
  "status": "healthy",
  "version": "2.1.0",
  "components": {
    "database": {
      "status": "connected",
      "latency_ms": 12
    },
    "cache": {
      "status": "connected",
      "hit_rate": 0.94
    },
    "queue": {
      "status": "connected",
      "pending_jobs": 42
    }
  },
  "metrics": {
    "requests_per_sec": 1250,
    "error_rate": 0.002,
    "p99_latency_ms": 145,
    "active_connections": 89
  },
  "uptime_seconds": 864000
}
      \end{lstlisting}
    \end{column}
  \end{columns}
\end{frame}

\begin{frame}[fragile]{Code in Blocks and Verbatim}
  \begin{exampleblock}{\scriptsize Configuration Example}
    \begin{lstlisting}[language=bash, numbers=none, basicstyle=\ttfamily\fontsize{5}{6}\selectfont]
# /etc/myapp/config.conf
server_address = "0.0.0.0"
server_port = 8443
tls_enabled = true
tls_cert = "/etc/pki/tls/certs/server.pem"
tls_key = "/etc/pki/tls/private/server-key.pem"
log_level = "info"
    \end{lstlisting}
  \end{exampleblock}

  \vspace{0.05cm}

  \begin{block}{\scriptsize Verbatim Text Example}
    {\fontsize{5}{6}\selectfont
    \begin{verbatim}
$ curl --cert client.pem --key client-key.pem \
    https://server:8443/api/v1/status
{"status":"ok","version":"2.1.0"}
    \end{verbatim}}
  \end{block}

  \vspace{0.05cm}

  \begin{alertblock}{\scriptsize Important Note}
    \tiny
    Use \texttt{[fragile]} frame option when including \texttt{lstlisting} or
    \texttt{verbatim} environments. Without it, LaTeX will fail to compile.
  \end{alertblock}
\end{frame}
